\chapter{Notación y símbolos}
\label{ch:notacion}

\begin{longtable}{@{}p{.16\textwidth}p{.30\textwidth}p{.42\textwidth}@{}}
\toprule
\textbf{Símbolo} & \textbf{Nombre} & \textbf{Significado en este libro}\\
\midrule
\endhead
\(\vect{v}=(u,v,w)\) & velocidad & campo euleriano \\
\(p\) & presión & pascal; en reposo, esfuerzo isótropo \\
\(\rho\) & densidad & kilogramo por metro cúbico \\
\(\mu\) & viscosidad dinámica & pascal segundo \\
\(\nu=\mu/\rho\) & viscosidad cinemática & metro cuadrado por segundo \\
\(\vect{g}=-g\vect{e}_z\) & gravedad & \(z\) hacia arriba, \(g>0\) \\
\(D/Dt\) & derivada material & \(\partial/\partial t+\vect{v}\cdot\nabla\) \\
\(\vect{\omega}=\nabla\times\vect{v}\) & vorticidad & radián por segundo en sentido físico \\
\(\psi\) & función de corriente & \(u=\partial\psi/\partial y\), \(v=-\partial\psi/\partial x\) \\
\(\varphi\) & potencial de velocidad & \(\vect{v}=\nabla\varphi\) si \(\vect{\omega}=\vect{0}\) \\
\(\mat{\sigma}\) & tensor de esfuerzos & \(\vect{t}=\mat{\sigma}\vect{n}\), Cauchy \\
\(\mat{D}\) & tensor de deformación & parte simétrica de \(\nabla\vect{v}\) \\
\(\mathrm{Re}\) & número de Reynolds & \(\rho V L/\mu\); \(L\) se declara \\
\(C_D\), \(C_L\) & coeficientes & fuerza partida por \(\tfrac12\rho V^2 A\) \\
\(f\) & factor de Darcy & \(\tau_w=f\rho V^2/8\) \\
\(\Omega\) & velocidad angular & del sistema o del cilindro, según el capítulo \\
\bottomrule
\end{longtable}
